\begin{explanationbox}
\textbf{深度解读：公式背后的逻辑}
\begin{itemize}
    \item \textbf{余弦定理的变体}：由数量积定义 $\vec{AB} \cdot \vec{AC} = |\vec{AB}||\vec{AC}|\cos A$。结合余弦定理 $\cos A = \frac{b^2+c^2-a^2}{2bc}$，消去分母中的边长项即可得到该公式。
    \item \textbf{极化恒等式的影子}：此公式与极化恒等式 $\vec{a}\cdot\vec{b} = \frac{1}{4}(|\vec{a}+\vec{b}|^2 - |\vec{a}-\vec{b}|^2)$ 的物理结构异曲同工，本质上是在处理三角形中的平行四边形模长关系。
    \item \textbf{应用场景}：在已知三边长度或涉及中线、投影计算的题目中，该公式可以绕过角度计算，直接进入代数运算逻辑。
\end{itemize}
\end{explanationbox}